% Generated draft: requires mathematical and source review.
\documentclass[11pt]{article}
\usepackage[margin=25mm]{geometry}
\usepackage{fontspec}
\setmainfont{Latin Modern Roman}
\usepackage{amsmath,amssymb,mathtools,unicode-math}
\setmathfont{Latin Modern Math}
\usepackage{enumitem,needspace,xurl,hyperref}
\setlength{\parindent}{0pt}
\setlength{\parskip}{5pt}
\setlength{\emergencystretch}{3em}
\newcommand{\N}{\mathbb{N}}
\newcommand{\Z}{\mathbb{Z}}
\newcommand{\Q}{\mathbb{Q}}
\newcommand{\R}{\mathbb{R}}
\newcommand{\C}{\mathbb{C}}
\newcommand{\F}{\mathbb{F}}
\newcommand{\A}{\mathbb{A}}
\newcommand{\PP}{\mathbb{P}}
\newcommand{\E}{\mathbb{E}}
\newcommand{\eps}{\varepsilon}
\newcommand{\dd}{\,\mathrm{d}}
\begin{document}
\section*{UnsolvedMath ID: 30005050}
\textbf{Source problem code:} catalogue entry
\subsection{Definitions and mathematical statement}
Let \(\zeta(s)\) denote the meromorphic continuation of \(\sum_{m=1}^{\infty}m^{-s}\), initially defined for \(\operatorname{Re}s>1\). Its nontrivial zeros are the zeros \(\rho=\beta+i\gamma\) with \(0<\beta<1\). Consider all such zeros with \(\gamma>0\), and list their ordinates in nondecreasing order:
\[
0<\gamma_1\leq\gamma_2\leq\cdots.
\]
For this draft, zeros are counted with their analytic multiplicities, so repeated ordinates are retained. This is a provisional interpretation of the source's non-strict ordering: the supplied record does not explicitly specify the multiplicity convention. No Riemann hypothesis or simplicity hypothesis is imposed in the supplied clean statement.

For \(n\geq1\), define the normalized consecutive gap by
\[
g_n=\frac{(\gamma_{n+1}-\gamma_n)\log\gamma_n}{2\pi},
\]
where \(\log\) is the natural logarithm. The conjecture asks whether both assertions hold:
\[
\liminf_{n\to\infty}g_n=0,
\qquad
\limsup_{n\to\infty}g_n=+\infty.
\]
Equivalently, for every \(\varepsilon>0\) and every integer \(N\geq1\), there exists \(n\geq N\) with \(g_n<\varepsilon\); and for every \(M>0\) and every integer \(N\geq1\), there exists \(n\geq N\) with \(g_n>M\).

The supplied excerpt motivates the normalization using the zero-counting asymptotic
\[
N(T)=\frac{T}{2\pi}\log\frac{T}{2\pi}-\frac{T}{2\pi}+O(\log T)
\qquad(T\to\infty),
\]
where \(N(T)\) counts nontrivial zeros with \(0<\operatorname{Im}\rho<T\), under the same multiplicity convention. The scale \(2\pi/\log T\) is an asymptotic mean-spacing scale. The first conjectured limit concerns gaps arbitrarily small relative to this scale, not gaps approaching this scale.

The original excerpt also contrasts these conjectures with a stated regular-spacing theorem for functions denoted \(\Xi_\lambda(t)\) when \(\lambda>0\). Those functions and the threshold \(T_\lambda\) are not defined in the supplied record, and their counting formula is typographically damaged. Consequently that comparison cannot be made into a precise auxiliary statement from the supplied evidence alone; it supplies no additional hypothesis for the conjecture above.
\subsection{Short English statement}
List the positive ordinates of all nontrivial Riemann-zeta zeros in order, provisionally retaining multiplicities. The conjecture asks whether consecutive gaps, divided by the asymptotic mean-spacing scale, become arbitrarily small and arbitrarily large along subsequences tending to infinity. The supplied formulation assumes neither the Riemann hypothesis nor simplicity of the zeros. Its multiplicity convention requires confirmation.
\subsection{Sources}
\begin{itemize}
\item Supplied local UnsolvedMath registry record, permanent ID \(30005050\), source code \texttt{catalogue entry}; registry location \texttt{data/unsolvedmath/problems.json}, with \texttt{catalog/problems.json} designated as the supplemental local registry. Title: \emph{Extreme Normalized Gaps between Riemann-Zeta Zeros}. No individual proposer is identified; the supplied proposed year is 2022.
\item Supplied report citation: \emph{The Laguerre-P\'olya Class and Combinatorics} (2023), Oberwolfach Reports, DOI \texttt{10.4171/owr/2022/13}, pp.\ 657--681. \url{https://doi.org/10.4171/owr/2022/13}. The literature-source label is ``Oberwolfach Report 13/2022.'' The supplied citation year 2023 and report designation 13/2022 are retained as given.
\item Supplied literature reference labelled ``Conrey--Ghosh--Gonek on zeta gaps'': \url{https://arxiv.org/abs/1003.0752}. The supplied record does not give a full bibliographic citation.
\item Supplied literature reference labelled ``Bui--Milinovich--Ng on zeta gaps'': \url{https://arxiv.org/abs/0910.2052}. The supplied record does not give a full bibliographic citation.
\item Literature assessment dated exactly \texttt{2026-08-22}: the supplied curation labels the problem open and describes the two extreme normalized-gap limits as conjectural, with cited primary work providing finite unconditional or conditional bounds. This is supplied dated curation, not an independently verified literature assessment. The background explicitly requests independent verification of time-sensitive status and cited claims.
\item The supplied statement status is \texttt{corrected\_verified}, accompanied by a claim that the clean statement was checked against the official Oberwolfach report. This is a supplied verification claim; no source visit or independent verification is asserted here.
\end{itemize}
\end{document}
